\section{Totality and open subgroups}

\subsection{The comparison theorem and its hypotheses}

\begin{theorem}[Proved: the general saturation theorem]
\label{B:saturation}
Assume (S) and (FA). For every covariant action $L:A_M\to\mathrm{Set}$
there are natural transformations
\[
 L\longrightarrow L^{\mathrm{inj}}\longrightarrow T_S,
 \qquad S=\operatorname{colim}_{\mathcal I_M}L^{\mathrm{inj}},
 \qquad T_S(B)=S^{G_B},
\]
each inducing a weak equivalence on categories of elements. The
$G$-set $S$ is smooth. Its orbits are in bijection with the connected
components of $\operatorname{El}(T_S)$. For every open $U\leq G$,
\[
 \operatorname{El}_{A_M}(T_{G/U})\simeq E_U
\]
as categories, where $E_U$ is the category in Lead 7. Thus $A$ is total
if and only if every $BE_U$ is weakly contractible; and $A\in\mathcal Q_n$
if and only if every $BE_U$ is $(n-1)$-connected. Neither implication
uses (FA$\varnothing$).
\end{theorem}
\begin{proof}
Amalgamation gives condition I. For $x,x'\in L(B)$ define $x\sim x'$
when $L(f)x=L(f)x'$ for some arrow $f$ out of $B$. Two witnesses can
be completed to a square, proving transitivity; the same argument
proves naturality under a map out of $B$. The quotients form
$L^{\mathrm{inj}}$, whose transition maps are injective by composition
of a witnessing arrow.

We first record the precise use of (FA). If $P$ is nonempty and
$\mathcal U$ is a nonempty full cosieve of $P/A_M$, choose $P\to Z$
in it. The functor $X\mapsto X\otimes_PZ$ comes with natural maps
from $X$ and from the constant object $Z$; its image stays in
$\mathcal U$ since it receives a map from $Z$. It contracts
$B\mathcal U$. The tagged finite amalgams may be made before passing
to $A_M$ and transported along an equivalence; no functorial choice
of embeddings into $M$ is involved.

For the quotient $q:L\to L^{\mathrm{inj}}$, fix $(B,[x])$ in the
target category of elements. The functor $B/A_M\to
(B,[x])\downarrow\operatorname{El}(q)$ sends $f:B\to C$ to
$(C,L(f)x)$. Its comma category at $((C,y),f)$ is the full cosieve
of $C/A_M$ consisting of $h$ with $L(h)y=L(hf)x$. It is nonempty
by the definition of the quotient and contractible by the preceding
paragraph. Two applications of Quillen's Theorem A prove that
$\operatorname{El}(q)$ is a weak equivalence.

For an injective action, the directed colimit $S$ has injective
structure maps $i_B:L(B)\to S$. The action
$g i_B(x)=i_{gB}(L(g|_B)x)$ is well-defined and smooth, exactly as
in Theorem~\ref{A:coinvariants}. Extension of a strong embedding
$f:B\to C$ to $g\in G$ defines $T_S(f)s=gs$, independently of that
extension. The maps $i_B$ give a natural injection $L\to T_S$.
At $(B,s)$ its comma category is the full cosieve of $B/A_M$ on
those $f:B\to C$ for which $T_S(f)s$ belongs to $i_C(L(C))$.
Injectivity makes the representing element unique. The cosieve is
nonempty: represent $s$ at some finite strong $D$ and take a finite
strong stage containing $B\cup D$. It is contractible by (FA),
and Theorem A proves the second weak equivalence.

For the last categorical equivalence choose a representative $g$
for each coset appearing in an object $(B,gU)$ of
$\operatorname{El}(T_{G/U})$, and send that object to $g^{-1}B$.
The fixed-coset condition is exactly $G_{g^{-1}B}\leq U$.
An arrow $f:B\to C$ extending to $h\in G$, with target coset $g'U$,
is sent to
\[
 (g'^{-1}hg)|_{g^{-1}B}:g^{-1}B\longrightarrow g'^{-1}C.
\]
Its extending element lies in $U$, and its restriction is independent
of the choice of $h$. These formulas give bijections on hom-sets;
every object of $E_U$ is the image of $(B,U)$, so this is an
equivalence. Different representative choices give naturally
isomorphic equivalences. Finally, the inclusion subposet of the
objects of $E_U$ is nonempty and directed, hence $E_U$ is connected.
This also proves that the orbit decomposition gives precisely the
components. Applying the all-action characterization of totality,
or its connectivity truncation, proves the criteria.
\end{proof}

\subsection{The two ordered equivalence structures}

\begin{lemma}[Proved: finite supports after adding the class sort]
\label{B:equivalence-supports}
Let $M$ be either the generic independently ordered equivalence
structure or the generic convexly ordered equivalence structure.
Write $J=M/E$. A finite closed support in the two-sorted structure
$(M,J,\pi)$ is a pair $(P,D)$, where $P\subset M$ and $D\subset J$
are finite and $\pi(P)\subseteq D$. Write $G_{P,D}$ for the subgroup
fixing all points of $P$ and all classes of $D$. Then
\[
 \langle G_{P,D},G_{P',D'}\rangle
   =G_{P\cap P',D\cap D'}.
\]
Every open subgroup $U$ has a unique least such finite support
$(P,D)$ with $G_{P,D}\leq U$, and is the inverse image of a subgroup
of the automorphism group of that finite support under its setwise
stabilizer. In the convex case this finite automorphism group is
trivial. In the independent-order case it permutes only
$D\setminus\pi(P)$, and is the full symmetric group there.
\end{lemma}
\begin{proof}
We spell out the generation argument to avoid importing an
elimination-of-imaginaries statement for a different language.
As usual, form the bipartite coset incidence graph of copies of the
two finite supports over their common support. To connect two
copies of the first support, first insert intermediate copies whose
new class names, apart from the common class names, avoid all
previously used classes. This is possible because there are
infinitely many generic classes, with a dense quotient order in
the convex case. Within each required class, points may be inserted
in every prescribed nonempty interval.

In the independent-order case the only further constraints are the
one linear order on the finitely many real points. Interpolate the
increasing tuples gap by gap over the common real points and
subdivide into finitely many overlapping-gap steps. At each step
a copy of the second support can be inserted so that both unions
have the prescribed type: its class equalities are already
consistent because noncommon names were chosen fresh, and its
point coordinates are chosen in the intersections of the required
gaps. In the convex case first perform the same overlapping-gap
construction for the finite class tuples in the quotient order,
then independently for the real points in each common class.
The finite extension property realizes every specified pattern.
Thus the incidence graph is connected. A finite path in the coset
graph is exactly a finite alternating product of the two
stabilizers, which proves equality as subgroups, not merely after
topological closure.

An open subgroup contains a finite real-point stabilizer and hence
one of the displayed support stabilizers. Choose a support of least
total size. The generation identity applied to that support and
its translate by any $u\in U$ makes their intersection a support
for $U$. Minimality forces setwise preservation of the chosen
support. It also proves uniqueness, by intersecting two minimal
supports. The kernel of the induced finite permutation action is
$G_{P,D}$, so $U$ is the inverse image of its image. The real-point
order forces every real point of a preserved finite support to be
fixed. In the convex case the quotient order also fixes every
class; in the independent case any permutation of the remaining
class markers extends by the dense-colour back-and-forth. These
are exactly the asserted finite automorphism groups.
\end{proof}

\begin{theorem}[Proved: all open-subgroup component types in A.4(2),(3)]
\label{B:equivalence-types}
For convexly ordered equivalence structures every open subgroup is
$G_{P,D}$ for a finite pair as in Lemma~\ref{B:equivalence-supports},
and every $BE_U$ is contractible. In the independently ordered case,
every open subgroup is determined by such a pair and a subgroup
$H\leq\operatorname{Sym}(D\setminus\pi(P))$, and
\[
 BE_U\simeq BH.
\]
In particular a single-class stabilizer gives a contractible
component in either case, while the stabilizer of an unordered pair
of classes in the independent-order case gives
$BC_2\simeq\mathbb{RP}^{\infty}$. More generally $BS_m$ occurs for
every $m$. The convex category is total, exact over every coefficient
ring, and Ramsey; its group is extremely amenable. Nevertheless
(Op) and (St) fail in that presentation.
\end{theorem}
\begin{proof}
The description of open subgroups follows from the lemma. A finite
real strong set $B$ satisfies $G_B\leq U$ exactly when it contains
$P$ and meets every class in $D$. Necessity for real points follows
from $\operatorname{Fix}(G_B)=B$. For class markers, $G_B$ moves
every class not represented in $B$ through infinitely many classes
of the appropriate quotient gap, so none can belong to the finite
preserved set $D$. Conversely these two conditions ensure
$G_B\leq G_{P,D}\leq U$.

First take $U=G_{P,D}$. In the abstract category of finite structures
containing the named real points $P$ and meeting the named classes
$D$, adjoin one new point to each distinguished class. In the
independent-order case place all new points above all existing
points, in one fixed order of the named classes. In the convex case
place the new point above all old points in its own class. This is
a functor under embeddings fixing $P$ and the named classes: extend
an embedding by the identity on each new formal point. There are
natural maps from the identity and from the fixed object consisting
of $P$ together with these new points. If both $P,D$ are empty,
use one new point in a new class. Thus this category is contractible.
Transporting tagged structures to the homogeneous limit gives an
equivalence with $E_{G_{P,D}}$, so that space is contractible.

For a general $U$, let $H=U/G_{P,D}$ be the finite group from the
lemma. The usual coset-labelled category, with objects $(B,h)$ and
arrows carrying the label by the induced class permutation, is a
regular covering of $E_U$ with deck group $H$. Choosing extensions
of the finitely many class permutations identifies this covering,
up to categorical equivalence, with $E_{G_{P,D}}$. It is therefore
contractible, proving $BE_U\simeq BH$.

For the convex class, finite linear orders of classes and finite
linear orders inside the classes have the Ramsey property by the
finite product Ramsey theorem; this is KPT Corollary~6.8, whose
statement and its reduction to ordered equivalence classes were
checked in the supplied source \cite{ABKPT2005}. The compact and
integral conclusions follow from Theorem~\ref{A:compact-ramsey}.
Alternatively, the just computed action components and (FA) give
totality directly. To see that (Op) fails, take the stabilizer of
one class without a named point; it is not a finite real-point
stabilizer. For two distinct real points of that same class,
the generated point stabilizers still preserve the class, whereas
their real supports have empty intersection; hence (St) fails.
\end{proof}

\subsection{Independent linear and cyclic orders}

\begin{proposition}[Proved: (St) and (Op) for the class of R10]
\label{B:cyclic-supports}
Let $C$ be the class of finite nonempty sets with an independent
linear and cyclic order. Its limit satisfies (St) and (Op), and
$\operatorname{Fix}(G_B)=B$ for every finite $B$. Every open subgroup
is uniquely $G_F$ for a finite set $F$. Under (FA), every $E_U$ is
contractible except possibly $E_G=C$, whose classifying space has
the nonzero second homology of R10. The group is not amenable and
is not extremely amenable.
\end{proposition}
\begin{proof}
The strong amalgamation and finite fixed-point assertion follow by
extending the two orders independently on a disjoint amalgam.
Functorial amalgamation over a nonempty base chooses a base point
as a cyclic cut and puts the first side before the second in every
linear gap and every cyclic gap.

For (St), let $D=B\cap C$ and use the incidence graph of copies of
$B$ and $C$ over $D$, with edges copies of their joint two-order
structure. Realize a finite configuration by two coordinates: a
real coordinate for the linear order and a circular coordinate for
the cyclic order. The configuration space of a tuple with a fixed
linear type and cyclic type is path connected. With $D$ nonempty,
each coordinate is confined to one of the gaps of $D$ and the
ordered-coordinate spaces are convex after cutting each gap. With
$D$ empty, the circular configuration space is a circle times an
open simplex, still path connected, independently of the convex
linear-coordinate space.

Join any two left configurations by a path in this product space.
Compactness of the parameter interval permits a finite subdivision
on which the corresponding required gaps overlap in both orders.
Perturb the intermediate positions to be pairwise distinct and
avoid the old finite configuration. For each neighbouring pair,
the vertices of a copy of $C$ can be placed in the intersections
of the required linear intervals and cyclic arcs, in their
prescribed orders. The two orders are independent, so this is a
consistent finite two-order structure. The extension property
realizes it in the limit over the old configuration. Hence every
two left vertices are connected by a finite alternating path,
proving (St). Finite intersections are admissible and finite
structures are rigid, so the minimal-support argument in
Proposition~\ref{B:intersection} proves (Op). Uniqueness follows
from the fixed-point identity. The component assertion follows
from Theorem~\ref{B:op-total}.

For nonamenability, complete the dense cyclic order of the limit
to a circle. Its automorphisms extend continuously to the circle,
and this gives a continuous $G$-action: fixing a sufficiently fine
finite cyclic mesh moves every point by an arbitrarily small amount.
Suppose there were an invariant probability measure. Individual
points of the countable dense orbit have measure zero, since that
orbit is infinite and transitive. Homogeneity of ordered pairs
implies that the positively directed arc from $x$ to $y$ has a
constant measure $t$ whenever $x<y$ in the independent linear order,
and measure $1-t$ when $y<x$.

Choose $a<b<c$ with cyclic order $(a,b,c)$. Their three consecutive
arcs partition the circle up to null endpoints and have masses
$t,t,1-t$; hence $t=0$. Choose instead $a<b<c$ with cyclic order
$(a,c,b)$, a second finite pattern present in the limit. Its
consecutive masses are $t,1-t,1-t$, forcing $t=1$. This is a
contradiction. Thus the action has no invariant probability
measure. The finite linear system is checked in the accompanying
certificate script. Non-extreme amenability follows immediately.
\end{proof}

\begin{remark}[Proved: scope of the free-amalgamation question]
\label{B:bridge-scope}
As literally stated for all strong classes in (S), B.4 is false:
the truncated linear-order example in Proposition~\ref{B:intersection}
is the free ordered expansion of the strong class of sets of size
zero or at least two. It has (St) but fails (Op). For the usual
all-substructure free-amalgamation classes, intersection closure is
automatic. The exact additional finite condition used by the
alternating argument is the following bridge condition: for each
two finite strong supports $B,C$ with intersection $D$, the
bipartite graph of their strong copies over $D$, whose edges are
strong copies of their joint structure, is connected. This condition
is equivalent to (St), since that graph is the coset graph
$G_D/G_B\sqcup G_D/G_C$.

One sufficient combinatorial way to verify it is closure under the
finite alternating cyclic amalgams used in T A.30, with the two
distinguished supports embedded strongly and with an order extending
their prescribed finite order patterns. For ordinary relational
free amalgamation this closure is proved in
Proposition~\ref{B:relational-bridge} below. For a
general strong class, free amalgamation along a tree does not by
itself assert that identifying the ends of that tree leaves an
admissible strong structure. The predimension positivity estimate
in Part D is precisely the missing verification in that application.
\end{remark}

\begin{proposition}[Proved: the fundamental group of an orbit component]
\label{B:fundamental}
Under (S), without (FA), every open $U\leq G$ satisfies
\[
 \pi_1(BE_U)\cong
 U\big/\langle G_B:B\text{ finite nonempty strong},\ G_B\leq U\rangle.
\]
The subgroup in the denominator is normal in $U$. Consequently
totality of $A$ implies that each open subgroup is generated by the
finite strong pointwise stabilizers which it contains.
\end{proposition}
\begin{proof}
Write $H$ for the denominator. Conjugation by $u\in U$ replaces
$G_B$ by $G_{uB}$, proving normality. An arrow $u|_B:B\to C$ of
$E_U$ defines the coset $uH$, independently of its extension: two
extensions differ on the right by $G_B$. This is a functor
$E_U\to B(U/H)$.

We compute the fundamental group using the edge-and-triangle
presentation of a category nerve. Its inclusion subposet is directed
and contractible, so all its edges may be collapsed in this
presentation. Every remaining edge is labelled by an element of
$U$. For $u\in U$, choose a finite strong $C$ containing
$B\cup uB$; the zigzag $B\xrightarrow{u|_B}C\leftarrow B$
represents $u$. Replacing $B$ by a larger support gives the same
loop by the commutative inclusion squares. Any two supports have a
common larger one, so the label does not depend on $B$. The triangles
give multiplication of labels, and a label in $G_B$ is the identity
edge on $B$, hence is killed. Conversely the functor just constructed
respects exactly these relations. This proves the displayed
isomorphism, including surjectivity and the absence of extra
relations.

The equivalence with the category of elements in the previous proof
did not use (FA). Therefore, if $A$ is total, $E_U$ is a connected
action component and has trivial fundamental group. The quotient
formula gives the conclusion.
\end{proof}

\begin{theorem}[Proved: the exact generalization of T A.33]
\label{B:op-total}
Under (S), (FA), and (Op), every action component is weakly
contractible or weakly equivalent to $BA$. Therefore $A$ is total
if and only if $BA$ is weakly contractible, and $A\in\mathcal Q_n$
if and only if $BA$ is $(n-1)$-connected. Adding (FA$\varnothing$)
implies totality. No separate rigidity hypothesis is required for
this theorem.
\end{theorem}
\begin{proof}
For every nonempty finite strong $B$,
$\operatorname{Fix}_M(G_B)=B$. Indeed, for $x\notin B$ choose a
finite strong $D$ containing $B\cup\{x\}$ and use the disjoint
amalgam in (FA) to make a second copy of $D$ over $B$. Embed this
finite strong extension into $M$ over the first copy and extend
the copy isomorphism by homogeneity. It fixes $B$ and moves $x$.
The same construction with arbitrarily many copies proves that
every $G_B$-orbit outside $B$ is infinite.

For finite nonempty strong $F$, the map
\[
 (G/G_F)^{G_B}\longrightarrow A_M(F,B),\qquad gG_F\longmapsto g|_F
\]
is well-defined and bijective. The fixed-coset condition says that
$gF\subseteq\operatorname{Fix}(G_B)=B$; since $gF$ is strong in
$M$, axiom (*) of (S) makes it strong in $B$. Homogeneity proves
surjectivity, and equality of restrictions is exactly equality of
cosets. The map is natural in $B$. Its category of elements is
$F/A_M$, with initial object $1_F$.

For $F=\varnothing$, $G/G_F$ is the singleton action, and its category
of elements is $A_M$ itself. By (Op), these are all orbit types; now
apply Theorem~\ref{B:saturation}. Finally, a functorial joint
embedding with any fixed nonempty $Z$ gives natural arrows from
the identity functor and the constant functor $Z$ to
$X\mapsto X\otimes_\varnothing Z$. It contracts $BA_M$.
\end{proof}

\begin{proposition}[Proved, with a necessary exception to Lead 7]
\label{B:finite-extension}
Assume (S) and (FA), and let $F$ be a \emph{nonempty} finite strong
set with $G_F\triangleleft U$ of finite index. Then
\[
 BE_U\simeq B(U/G_F).
\]
If $F=\varnothing$, the conclusion must instead be $BE_G\simeq BA$.
Thus the finite-extension assertion of Lead 7 is false as written
when it includes the empty support and assumes only (FA).
\end{proposition}
\begin{proof}
Normality and $\operatorname{Fix}(G_F)=F$ imply that $U$ preserves
$F$ setwise. If $G_B\leq U$, the set $F$ is a finite $G_B$-set.
The infinite-orbit assertion in the preceding proof gives $F\subseteq B$.
Conversely, $F\subseteq B$ implies $G_B\leq G_F\leq U$. Thus the
objects of $E_U$ are exactly the finite strong supersets of $F$.

Put $\Gamma=U/G_F$. Form a category $\widetilde E$ with objects
$(B,\gamma)$, $B\in E_U$, $\gamma\in\Gamma$, and arrows
$u|_B:(B,\gamma)\to(C,\gamma')$ when
$\gamma'=\bar u\gamma$. The object $(F,1)$ is initial: its unique
arrow to $(B,\gamma)$ is the restriction to $F$ of any representative
of $\gamma$. Hence $B\widetilde E$ is contractible. The free deck
action $\gamma\mapsto\gamma\eta$ has quotient $E_U$, and its nerve
is a simplicial covering. This is a universal cover with deck group
$\Gamma$, proving the formula.

For empty $F$ we have $G_F=G$, necessarily $U=G$, and $E_G=A_M$ by
definition. The cocycle class $A_2(\mathbb F_2)$ satisfies (S), (FA),
and (Op), but has classifying space $K(\mathbb F_2,2)$; it is the
explicit counterexample to the omitted empty-support qualification.
\end{proof}

\begin{corollary}[Proved: the WEI criterion and its empty-base scope]
\label{B:wei}
Under (S), (FA), (FA$\varnothing$), and (WEI), quasi-filteredness,
totality, and (Op) are equivalent. Without (FA$\varnothing$), this
three-way equivalence is false: $A_2(\mathbb F_2)$ is quasi-filtered,
satisfies (FA), (WEI), and (Op), and is not total.
\end{corollary}
\begin{proof}
A quasi-filtered category has simply connected action components,
by the condition-I/IV characterization rechecked in DF A.2.
For open $U$, choose the support supplied by (WEI). If it is
nonempty, Proposition~\ref{B:finite-extension} identifies the
fundamental group of $E_U$ with $U/G_F$, so it must be trivial.
If it is empty then $U=G$ already. Thus (Op) follows. Theorem
\ref{B:op-total} gives totality, and totality implies
quasi-filteredness. For the last assertion the cocycle class has
all nontrivial orbit components contractible and its singleton
component $K(\mathbb F_2,2)$, which is simply connected but not
contractible. The verified action characterization gives precisely
the stated properties.
\end{proof}



\subsection{The orbit-category dictionary and its variance}

\begin{theorem}[Proved with corrections: Lead 12]\label{B:orbit-dictionary}
Assume (S) and (SA), let $R$ be a commutative ring, and let $\mathcal P$ be the conjugacy-invariant
collection of subgroups $G_B$ for nonempty finite strong $B$.
Let $\mathcal O_{\mathcal P}$ be the full category of smooth
transitive $G$-sets $G/G_B$ and $G$-equivariant maps. Then
\[
 A^{\mathrm{op}}\simeq\mathcal O_{\mathcal P}.
\]
Thus a covariant $R$-module diagram on $A$ is a contravariant
(right) $\mathcal O_{\mathcal P}$-module. With $\underline R$
denoting the constant \emph{covariant} (left) module,
\[
 \operatorname{colim}_A V
 \cong V\otimes_{\mathcal O_{\mathcal P}}\underline R,\qquad
 T_i(V)\cong
 \operatorname{Tor}^{\mathcal O_{\mathcal P}}_i
       (V,\underline R).
\]
Exactness is flatness of this constant left module. This is an
orbit-category homology statement, but it is not, without changing
variance and checking the collection, the usual Bredon
homological-dimension-zero assertion about the constant right
module.

There is a contractible simplicial $G$-space $X_A$, with open
isotropy groups from $\mathcal P$, whose orbit space is naturally
$B(A^{\mathrm{op}})$, hence homeomorphic to $BA$. Its $H$-fixed
space is the nerve of the category of finite strong embeddings
whose image is fixed pointwise by $H$. It is a classifying space
for a subgroup family only when the corresponding fixed spaces
are contractible for every member of that family; this extra
condition is not supplied merely by specifying a collection.
\end{theorem}
\begin{proof}
Strong amalgamation gives $\operatorname{Fix}_M(G_C)=C$ for
nonempty strong $C$. Indeed, if $x\notin C$, choose a finite
strong $D$ containing $C\cup\{x\}$, amalgamate two copies of
$D$ strongly over $C$, realize the amalgam over $D$ in the
limit, and extend the resulting partial isomorphism. This
moves $x$ while fixing $C$.

An embedding $f:B\to C$ extends to $g\in G$. Associate to it
\[
 G/G_C\longrightarrow G/G_B,\qquad xG_C\longmapsto xgG_B.
\]
The subgroup inclusion $G_C\leq gG_Bg^{-1}$ makes it
well-defined. Two extensions of $f$ differ on the right by
$G_B$, so the map is independent of the extension. Conversely a
$G$-map between these orbits is determined by a coset $gG_B$
fixed by $G_C$. The fixed-point identity gives $gB\subseteq C$.
The subset $gB$ is strong in $M$ and therefore strong in $C$
by (*), so $g|_B$ is an arrow. Equality of restrictions is
exactly equality of the cosets. For composable embeddings
extended by $g$ and $h$, the composite is extended by $hg$;
the displayed orbit maps compose by the same formula in reverse
order. This proves fullness, faithfulness, and variance.

The tensor product is the direct sum of $V(G/G_B)$ modulo
the relations $V(\phi)v\otimes1=v\otimes1$ for orbit arrows
$\phi$. These are precisely the colimit relations on $A$.
Taking the standard representable bar resolution proves the
derived identity. In the standard Bredon convention, the
cellular fixed-point chains are contravariant and resolve a
constant \emph{right} module; covariant coefficients are then
tensored against them. L\"uck, Section~5.2, and Nucinkis,
Section~2, use exactly those variances
\cite{ABLueck2005,ABNucinkis2004}. Replacing a collection by its
downward closure changes its orbit category, so no equality of
these derived functors follows without an additional comparison.

For the geometric assertion, consider the covariant orbit
functor $\mathcal O_{\mathcal P}\to\mathrm{Set}$ sending $G/G_B$
to its underlying set, and form its category of elements
$\mathcal D$. Concretely its objects are strong embeddings
$e:B\to M$, and a morphism from $(C,e_C)$ to $(B,e_B)$ is an
embedding $f:B\to C$ with $e_Cf=e_B$. Sending an embedding to its
image gives an equivalence from $\mathcal D$ to the opposite
of the directed poset of finite strong subsets of $M$. Thus
$X_A=|N\mathcal D|$ is contractible. The action of $G$ on the
embeddings is simplicial and has open point stabilizers. An
$n$-simplex consists of a chain of orbit arrows and a point in
the initial orbit; $G$ is transitive on that orbit. Hence
$(N\mathcal D)/G=N\mathcal O_{\mathcal P}$, degree by degree.
Realization commutes with this quotient. Fixed simplices have
their embedded images fixed pointwise by $H$, giving the
asserted description of $X_A^H$.

For ordinary relational ages with all finite subsets admissible,
the finite subsets of $\operatorname{Fix}_M(H)$ form a directed
poset whenever nonempty. In that case the fixed-point
contractibility criterion does identify this space with the
classifying space for the downward-closed family generated by
the point stabilizers. For general strong classes the analogous
directedness must be checked. L\"uck's Theorem~1.9 is the
fixed-point characterization; its published survey assumes
locally compact groups, so here the needed open-isotropy
simplicial case has been proved directly. No local compactness
of a Fra\"iss\'e automorphism group is being assumed.
\end{proof}

\begin{proposition}[Proved: an explicit model for every $E_U$]
\label{B:orbit-model}
In setting (S), for each open $U\leq G$ form a category
$\mathcal D_U$ with objects $(B,e)$, where $B$ is an object of
$E_U$ and $e:B\to M$ is the restriction of an element of $U$.
An arrow $(C,e_C)\to(B,e_B)$ is an $E_U$-arrow $f:B\to C$
with $e_Cf=e_B$. Then $\mathcal D_U$ is contractible,
$U$ acts on it simplicially by postcomposition, and
\[
 BE_U\ \cong\ |N\mathcal D_U|/U
\]
after the canonical reversal of nerve coordinates. Its isotropy
groups are exactly the pointwise stabilizers occurring at the
finite images. This gives a model in every homotopy degree,
including for subgroups not finite extensions of pointwise
stabilizers. Under (SA) and (Rig), the images identify
$\mathcal D_U$ equivariantly, up to equivalence, with the
opposite of the poset of finite strong subsets $B$ with
$G_B\leq U$. That poset is identified with
$\{V\in\mathcal P:V\leq U\}$ ordered by reverse subgroup
inclusion, with $U$ acting by conjugation.
\end{proposition}
\begin{proof}
The image of every object is a finite strong $C$ with
$G_C\leq U$. Conversely each such image is obtained by taking
$e$ to be the identity. Between two objects there is a unique
arrow exactly when the second image is contained in the first;
it is the required restriction of an element of $U$.
Consequently the image functor is an equivalence with the
opposite directed poset of these finite strong subsets.
Directedness follows by placing any two in one larger finite
strong subset; its pointwise stabilizer is smaller. The nerve
is therefore contractible.

For a fixed chain of $E_U^{\mathrm{op}}$-arrows, the possible
initial embeddings form the transitive set $U/G_B$.
Quotienting its nerve by $U$ leaves exactly that chain, proving
the displayed orbit-space formula. The stabilizer of an
embedded object is its pointwise stabilizer, and fixes the
whole simplex whenever it fixes the initial embedded object.
Under (SA), the fixed-point identity makes
$B\mapsto G_B$ injective, and $G_C\leq G_B$ is equivalent to
$B\subseteq C$. Under (Rig) the setwise and pointwise
stabilizers of $B$ coincide, so this identification also respects
all isotropy groups and the conjugation action. This proves
the last description. No assertion that an orbit of a
contractible space is contractible is used.
\end{proof}

\begin{proposition}[Proved: a stronger necessary group condition]
\label{B:no-open-quotients}
In setting (S), if $A$ is total, no open subgroup $U\leq G$
has a proper open normal subgroup. In particular (Fin) holds.
This conclusion requires neither (FA) nor (SA).
\end{proposition}
\begin{proof}
Suppose $V\lhd U$ is proper and open. Inside $E_U$ take the
full subcategory $D$ on the objects with $G_B\leq V$. It is
nonempty because finite strong point stabilizers form an open
basis. It is connected because its inclusion poset is
directed. It is a full cosieve: for an arrow $u|_B:B\to C$,
\[
 G_C\leq uG_Bu^{-1}\leq V,
\]
using normality in $U$.

Totality passes to an action component by composing actions,
and to a nonempty connected full cosieve by extending actions
by the empty set. Since $E_U$ is equivalent to
$\operatorname{El}(T_{G/U})$, totality would make $D$ weakly
contractible. But there is a functor
\[
 D\longrightarrow B(U/V),\qquad (u|_B)\longmapsto uV.
\]
It is well-defined: two extensions of the same restriction
differ on the right by $G_B\leq V$. It induces a surjection on
fundamental groups. To see this, fix $B$ in $D$. For any $u\in U$,
choose a finite strong $C$ containing $B\cup uB$. The two
arrows $B\to C$, inclusion and $u|_B$, form a loop in the
groupoid completion whose image is $uV$ or its inverse.
Thus $\pi_1(D)$ surjects onto the nontrivial group $U/V$,
a contradiction. Finally, a proper finite-index open subgroup
has a proper open normal core, so (Fin) follows.
\end{proof}

\subsection{Step-by-step audit of the cocycle examples}

\begin{proposition}[Known and rechecked: DF A.4--A.8]
\label{B:cocycle-audit}
For countable abelian $\Lambda$ and $r\geq2$, the class of finite
ordered $\Lambda$-valued $r$-cocycles, with all induced embeddings,
satisfies (S), (Rig), (SA), (FA), (St), and (Op).
Every action component is weakly contractible or its projection
to $A_r(\Lambda)$ is a weak equivalence, and
\[
 BA_r(\Lambda)\simeq K(\Lambda,r).
\]
Consequently
$A_n(\mathbb F_2)\in\mathcal Q_n\setminus\mathcal Q_{n+1}$
for every $n\geq2$. For its group, $U=G$ is an open subgroup for
which $E_U$ is $(n-1)$-connected but not $n$-connected.
The parameter condition $r\geq2$ is essential to the stabilizer
proof.
\end{proposition}
\begin{proof}
We record the audit with the actual extension mechanisms. On an
ordered finite set $S$, use alternating normalized simplicial
cochains, zero on tuples with repeated entries. In positive
degrees, adjoining a specified vertex $p$ gives the contraction
\[
 (s_pc)(v_0,\ldots,v_{r-1})=c(p,v_0,\ldots,v_{r-1}),
 \qquad ds_p+s_pd=1.
\]
If an $r$-cocycle is given on a finite union $K$ of simplices and
$H^r(K;\Lambda)=0$, write it as $db$, extend $b$ by zero to the
full simplex, and take its coboundary. This is a genuine cocycle
extension, not merely an extension of the separate values.
For two simplices meeting in a simplex, possibly empty, the
required cohomology vanishes. This proves hereditary closure,
joint embedding, and strong amalgamation. Countability follows
from countability of $\Lambda$; the rich-chain Fra\"iss\'e
construction gives (S). The reference order gives rigidity.

For (FA), fix a nonempty base $P$ and choose $p=\min P$.
On $X$ and $Z$ take the primitives $s_pc_X$ and $s_pc_Z$.
They agree on $P$. On the tagged union over $P$ retain these
values on either side and put zero on every mixed tuple.
Its coboundary is the amalgamated cocycle. Put first-side points
before second-side points in every gap of $P$. A map on the
first side fixing $P$ preserves the three tuple cases, and
pullback commutes with $s_p$; hence it induces an increasing
cocycle embedding of the amalgams, compatibly with composition.
This verifies the functoriality used by the saturation theorem.

For the fixed-point assertion, strongly amalgamate arbitrarily
many copies of $B\cup\{x\}$ over $B$, when $x\notin B$, and use
homogeneity. Thus $\operatorname{Fix}(G_B)=B$, and indeed all
other $G_B$-orbits on points are infinite. Put $D=B\cap C$.
The bipartite graph whose vertices are copies of $B$ and $C$
over $D$, and whose edges are copies of $B\cup C$, is the coset
graph of $G_D/G_B$ and $G_D/G_C$. Its connectedness is exactly
$\langle G_B,G_C\rangle=G_D$, with no closure of the generated
subgroup taken.

Consider two copies $B_0,B_1$ disjoint off $D$ whose corresponding
reference-order gaps overlap wherever points of $C\setminus D$
must be inserted. Insert a fresh copy $C'$ in those overlaps.
Prescribe the existing cocycle on $B_0\cup B_1$, and the required
copies of the cocycle of $B\cup C$ on $B_0\cup C'$ and
$B_1\cup C'$. These prescriptions agree on their pairwise
intersections, which are $B_0,B_1,C'$. Their triple intersection
is $D$. When $D\ne\varnothing$, the nerve of this three-simplex
cover is a filled triangle. When $D=\varnothing$, it is a circle.
The intersections are simplices, so elementary Mayer--Vietoris
(or the finite good-cover nerve lemma) gives
$H^r(K;\Lambda)=0$ in either case for $r\geq2$.
The extension construction above realizes $C'$ as a common
neighbor in the coset graph. This is the exact place at which
$r\geq2$ enters: $H^1(S^1;\Lambda)$ need not vanish.

For arbitrary $B_0,B_1$, in each interval of the ordered set cut
out by $D$, interpolate linearly between their increasing tuples
in real coordinates. The tuple space is convex; a sufficiently
fine finite mesh, with small perturbations of its interior
tuples, gives a chain of fresh tuples with overlapping neighboring
gaps. Prescribe the cocycle of $B$ separately on the intermediate
copies, and the actual cocycle on the old union $B_0\cup B_1$.
The resulting pieces meet only in $D$; their union is contractible
if $D$ is nonempty and a disjoint union of simplices otherwise,
so its positive-degree cohomology vanishes. Extend the cocycle,
then embed into $M$ over the old union. Consecutive copies have
a common neighbor by the preceding paragraph. Thus the coset
graph is connected and (St) holds. All finite intersections
remain members of this class, so the minimal-support argument
of Proposition~\ref{B:intersection} applies: for an open $U$,
choose $F$ of smallest cardinality with $G_F\leq U$; then
$G_{F\cap uF}\leq U$ for $u\in U$ forces $uF=F$, and rigidity
forces $u$ to fix $F$ pointwise. Therefore $U=G_F$.

For the homotopy type, let
$K_m=Z^r(\Delta^{[m]};\Lambda)$. Pullback, with repeated tuples
zero, makes $K_\bullet$ a simplicial abelian group. Its normalized
Moore complex has $\Lambda$ in degree $r$ and zero elsewhere:
below $r$ there are no relevant tuples; in degree $r$ the value
on the full tuple is free; in degree $r+1$ the only candidate
tuple surviving the prescribed face kernels is killed by the
cocycle equation; in higher degrees no tuple can contain all the
required vertices. The normalized-complex comparison, specifically
\cite[Theorem III.2.4 and Corollary III.2.7]{ABGoerssJardine},
therefore gives $|K|\simeq K(\Lambda,r)$. The category is
$\Delta_{\mathrm{inj}}/K$, with degenerate simplices retained.
The inclusion into $\Delta/K$ is a weak equivalence by
Theorem~A: the comma at an $m$-simplex consists of nondecreasing
maps into $[m]$ with injective maps of domains, and prepending
a vertex labelled $0$ gives natural arrows from the identity
and the constant one-vertex object to one endofunctor.
The simplex-category realization is $|K|$. This checks the
variance and the degeneracy issue in DF A.4.

Theorem~\ref{B:saturation} now applies to all actions, including
noninjective and infinite ones. By (Op), their saturated orbit
components with support $F\ne\varnothing$ are represented by
$F/A$ and have an initial object; the empty support gives the
constant singleton action with element category $A$ itself.
The comparison maps in Theorem~\ref{B:saturation} are over $A$,
so in the latter case the original component's projection is
a weak equivalence. This is the full conclusion of DF A.7.

For completeness, the integral cycle of DF A.8 can be checked
without relying on a mod-two rank computation. Let $t$ be the
$n$-simplex with top cocycle value $1\in\mathbb F_2$, and $0$
the zero simplex. Their proper faces coincide, hence $t-0$ is
an integral cycle. Its barycentric subdivision is the alternating
sum, over permutations $\sigma$ of $0,\ldots,n$, of the flags
\[
 \{\sigma(0)\}\subset\{\sigma(0),\sigma(1)\}\subset
 \cdots\subset\{0,\ldots,n\},
\]
with coefficient $\operatorname{sgn}(\sigma)$, first in $t$ and
then subtracted in $0$. Interior boundary faces cancel in pairs;
the top-deletion faces cancel between the two simplices.
The last-vertex cocycle over $\mathbb F_2$ evaluates to one:
only the increasing permutation has strictly increasing
successive maxima. Therefore this is not an integral boundary.
The Moore calculation and Hurewicz identify its homology group
as $\mathbb F_2$, so the class has order exactly two.
Every action component is $(n-1)$-connected, but the constant
singleton action has this nonzero $H_n$, giving the asserted
strict tower and the subgroup $U=G$.
\end{proof}

\subsection{A missing intersection hypothesis in the support argument}

\begin{proposition}[Refuted as literally stated; repaired with (Int)]
\label{B:intersection}
The assertion ``(Rig) and (St) imply (Op)'' in setting (S) requires
an additional intersection condition. Let (Int) mean that an
intersection of finite strong subsets of $M$ is again an admissible
finite strong subset, with the empty subset allowed. Under (Int),
(Rig), and (St), the assertion is true. Under the literal
pointwise-stabilizer reading of (St), it is false without (Int).
\end{proposition}
\begin{proof}
With (Int), choose $B$ of least cardinality among strong supports
with $G_B\leq U$. For $u\in U$, both $G_B$ and $G_{uB}$ lie in $U$,
so (St) gives $G_{B\cap uB}\leq U$. Minimality and (Int) force
$B=uB$. Rigidity then says that $u$ fixes $B$ pointwise, so
$U=G_B$. If the minimum is zero, the same conclusion is $U=G$.

For a counterexample, let $\mathcal K$ consist of finite linear
orders of cardinality zero or at least two, with all increasing
embeddings between its objects. This is a strong class in exactly
the sense of (S): it has countably many types, joint embeddings,
amalgamation, and axiom (*). Its rich limit is $(\mathbb Q,<)$.
All its nonempty objects are rigid, and the usual gap-wise ordered
amalgams prove (FA) and (FA$\varnothing$). Its strong subsets are
the finite subsets of size at least two, together with the empty
subset.

For the rational order, $\langle G_B,G_C\rangle=G_{B\cap C}$ for
all finite $B,C$. Here is a direct check. The coset incidence graph
has left vertices copies of $B$ over $D=B\cap C$, right vertices
copies of $C$ over $D$, and edges copies of $B\cup C$. Two left
tuples whose corresponding gaps overlap have a common right
neighbour. In every gap of $D$, interpolate between the two
increasing real-coordinate tuples, subdivide sufficiently finely,
and perturb intermediate coordinates to avoid coincidences. The
successive tuples have overlapping corresponding gaps. Density and
homogeneity of $\mathbb Q$ realize this finite order pattern. Thus
the incidence graph is connected, which is precisely the displayed
generation identity. This proves (St) even when $B\cap C$ has
one point.

The open subgroup $G_{\{x\}}$ is not $G_F$ for any allowed strong
$F$: $F=\varnothing$ gives the whole group, while for nonempty $F$
the identity $\operatorname{Fix}(G_F)=F$ forces $F=\{x\}$, an
excluded object. Hence (Op) fails. The category is nevertheless
Ramsey: an ordinary finite Ramsey target for objects of size at
least two still lies in the class. It is therefore exact and total.
If one intends the notation (St) itself to include intersection
closure, that is an additional hypothesis, and the first paragraph
is the corrected statement rather than a counterexample to that
stronger convention.
\end{proof}

% Standalone proof note for integration into Part B; AI-produced GPT-6 Astra.
\subsection{The alternating bridge for independently ordered posets}
\begin{proposition}[Proved: finite bridge for induced poset embeddings]
\label{B:poset-bridge}
Let $B,C$ be induced subposets of a finite poset $D$, with an independent
linear order on all structures, and put $F=B\cap C$. Let
$d_0,d_1:D\to E$ be increasing induced-poset embeddings agreeing on $F$.
There is a finite independently ordered poset $Z$, an increasing induced
embedding $E\to Z$, and increasing induced embeddings
$a_0,\ldots,a_{2N}:D\to Z$ with the prescribed endpoints, such that
$a_{2j}|_B=a_{2j+1}|_B$ and
$a_{2j+1}|_C=a_{2j+2}|_C$.
\end{proposition}
\begin{proof}
Use the strict set quotient of one $E$-copy and $2N-1$ copies of $D$
around a circle, identifying alternate overlaps over $B$ and $C$ and
using $d_0,d_1$ at the ends. Write $s=2N$ and require $s>|D|+4$.
As in the explicit fiber calculation in T~A.30, $F$-points have full
visibility and every other point has a proper visibility interval of
at most three consecutive factors. Each factor injects, and adjacent
factor intersections are precisely the prescribed induced copies of
$B$ or $C$. Give the quotient the reflexive transitive closure of all
copied poset relations. We must prove antisymmetry and inducedness.

We use the following elementary path argument, including its reduction.
A path in this generated relation can be described as a succession of
relations inside factors, crossing between consecutive factors at points
of their overlap. If two successive relations were originally recorded
in nonconsecutive factors, insert identity segments along the common
point's visibility interval; for a point outside $F$ this has length at
most two. Split the path at every point of $F$, so every resulting piece
has no $F$-point in its interior. Delete any immediate factor backtrack
$i,j,i$: the endpoints of the segment inside factor $j$ lie in its
induced overlap with factor $i$, so the same relation holds in factor $i$.
This strictly shortens the factor word and does not introduce an $F$
point. Thus a reduced word either has one factor or moves monotonically
along the factor circle.

No such reduced piece can wind around the circle. In fact, a winding
piece crosses the $s-1$ consecutive interior $D$-factors. Inside each
interior factor between two boundary crossings, its boundary labels
belong alternately to $B$ and $C$ and form a comparable pair in $D$.
Equality of these labels would put the boundary point in $F$.
Except possibly at the two endpoints of the piece, the comparison is
therefore strict. All these labels form an increasing chain in the
same finite poset $D$, of length at least $s-3>|D|$, impossible.
The possible three-index visibility through $E$ affects only the
$E$-segment, not this consecutive chain of interior $D$-factors.

Consequently each $F$-free piece misses an overlap. Cut there. Its
entire path lifts to the linear free amalgam of the factors; no extra
identification is crossed by the path. A linear free amalgam of posets
along induced subposets is a poset and contains every factor induced:
for a two-factor amalgam, a cross relation holds exactly when it is
mediated by a base point, and any path returning to one factor contracts
through the induced base; induction proves the linear assertion.

If both endpoints of the original path lie in one factor $T$, and the
path has no $F$-point, choose the endpoint occurrences in $T$ when
forming the factor word (identity transitions within their visibility
intervals suffice). The cut argument proves its relation already in
$T$. If the path does meet $F$, its first and last pieces similarly
give the relations from its first endpoint to a point of $F$, and from
a point of $F$ to its last endpoint, inside $T$. Each intervening piece
has endpoints in $F$ and yields a relation already in $F$, since $F$
is contained induced in every factor of the cut linear amalgam.
Composition proves the entire relation in $T$. Thus every factor is
induced in the generated preorder.

For antisymmetry, a directed cycle disjoint from $F$ reduces to a
closed factor walk. If it does not wind it lifts to a linear free
amalgam and is constant; winding was just excluded. A directed cycle
meeting $F$ splits into pieces whose endpoint relations lie in $F$.
Antisymmetry of $F$ makes all those endpoints equal, and antisymmetry
of each corresponding cut linear amalgam then makes each piece
constant. Hence the quotient preorder is a poset.

The independent linear order is obtained by the interpolation argument
of T~A.30. Embed $E$ increasingly in $\mathbb R$ and interpolate
$p_t=(1-t)d_0+td_1$ on $D$, constant on $F$. For sufficiently large
$N$, alternately hold the $B$-coordinates fixed and move the other
coordinates from $p_{j/N}$ to $p_{(j+1)/N}$; every intermediate map
remains strictly increasing. These maps respect the quotient equations,
so ordering the quotient by their real coordinates and breaking ties
arbitrarily preserves all factor linear orders. There is no requirement
that this independent order extend the poset relation. This gives the
claimed finite bridge.
\end{proof}

\begin{corollary}[Proved: (St) and (Op) for the independently ordered generic poset]\label{B:poset-op}
The category of finite posets with an independent linear order and all
increasing induced-poset embeddings satisfies (SA), (FA), (FA$\varnothing$),
(St), and (Op). Its category is total by the all-action theorem in Part B.
\end{corollary}
\begin{proof}
The free-poset amalgam $X\cup_P Z$ has cross relations mediated through
$P$. Its two factors meet exactly in $P$. Given an induced embedding
$u:X\to Y$ over $P$, the map $u\cup\operatorname{id}_Z$ is induced:
a cross comparability $u(x)\leq z$ in the target is equivalent to
$u(x)\leq p\leq z$ for some $p\in P$, which is equivalent to
$x\leq p\leq z$ since $u$ reflects order. Similarly for $z\leq u(x)$.
Choose the $X$-side before the $Z$-side in each gap of the independent
linear order on $P$. This is functorial, and over the empty base it is
disjoint union with concatenated linear order.

For the limit $M$, any finite subsets are strong. Given
$g\in G_{B\cap C}$, take finite $D$ containing $B\cup C$ and finite
$E$ containing $D\cup gD$, and apply the bridge to $1_D,g|_D$.
Embed $Z$ in $M$ over $E$ and extend its $D$-copies by homogeneity.
Consecutive automorphisms differ alternately by elements of $G_B$
and $G_C$, giving $g\in\langle G_B,G_C\rangle$. The other inclusion
is immediate. Minimal finite supports and rigidity of a finite
linear order then give (Op), with the empty support allowed, exactly
as in the proof of Theorem D.2.
\end{proof}

\begin{proposition}[Proved: the ordinary relational free-amalgamation case of B.4]
\label{B:relational-bridge}
Let $\mathcal K$ be a hereditary class of finite structures in a relational
language, with all induced embeddings, joint embedding, and free
amalgamation. Expand each member by every possible independent linear
order. The expanded class has (St) and (Op). More precisely, every
alternating bridge of Proposition~\ref{B:poset-bridge} exists in this
expanded class.
\end{proposition}
\begin{proof}
Use the same strict set quotient of $E$ and $s-1$ interior copies of $D$,
where $s$ is even and at least eight, and interpret each relation as the
union of its interpretations in the factors. The overlap identifications
use the two prescribed embeddings $d_0,d_1$ at the ends. An element of
$F=B\cap C$ has full visibility around the factor circle. Every other
element has a visibility interval of at most three consecutive factors:
it cannot cross both boundaries of an interior $D$-factor unless its
label belongs to $B\cap C$, and only the distinguished $E$-factor can
identify different endpoint labels. This also proves that the factors
inject into the quotient.

Each factor is induced. To see this, suppose that a tuple in one factor
$T$ acquires a relation from another factor $T'$. If all its entries are
in $F$, the two interpretations already agree on $F$. Otherwise choose
one entry outside $F$. Its visibility interval has at most three
vertices of the factor circle and contains $T,T'$. Every other entry
has full visibility or a short interval containing the same endpoints.
For a circle of length at least eight, all such short intervals contain
the same short arc between those endpoints. Thus the whole tuple lies
in every overlap along that arc. Since adjacent factors agree on their
induced overlap, its relation propagates back to $T$.

It remains to show that the quotient belongs to $\mathcal K$. Any
finite induced structure minimal outside a hereditary free-amalgamation
class is irreducible: any two distinct vertices occur together in a
relation tuple. Indeed, if $a,b$ did not occur together, the structure
would be the free amalgam of its deletions of $a$ and $b$ over their
intersection, and both deletions belong to the class. This contradicts
free amalgamation.

Every irreducible induced substructure $Y$ of our quotient lies in one
factor. For any two vertices of $Y$, a relation witnessing
irreducibility comes from a factor, so their visibility intervals
intersect. Ignore the intervals that are the whole circle. If any
proper interval remains, choose one, $I$, of length at most three.
Every other proper interval meets $I$ and has length at most three,
so their union is contained in $I$ enlarged by two vertices at each
end, a set of at most seven factor vertices. Cut the circle outside
that set. The remaining intervals are intervals of a line and are
pairwise intersecting; choosing the largest left endpoint shows that
they have a common vertex. The full-visibility intervals contain it
as well. If there were no proper intervals, $Y\subseteq F$ already.
Thus $Y$ is induced in one original factor and belongs to
$\mathcal K$. A hypothetical smallest induced substructure of the
quotient outside $\mathcal K$ would be irreducible, which is impossible.
Therefore the quotient is in $\mathcal K$.

For completeness, the order extension is independent of the relations.
Embed the finite order on $E$ into $\mathbb R$ and put
$p_t(x)=(1-t)d_0(x)+td_1(x)$ for $x\in D$. These maps are increasing
and constant on $F$. Subdivide $[0,1]$ sufficiently finely. Between
$p_t$ and $p_{t+h}$ use the map equal to $p_t$ on $B$ and to
$p_{t+h}$ on $D\setminus B$; it is increasing because the uniformly
positive gaps in the finite family $p_t$ dominate the sufficiently
small coordinate changes. It agrees with $p_{t+h}$ on $C$, since
$B\cap C=F$. Alternating these maps with the $p_t$ gives all the
required overlap equations. Their coordinates descend to the quotient;
break ties between distinct points in any order. No factor has a tie,
so its prescribed order is preserved.

All embeddings are induced and therefore strong. The coset-graph
argument in Corollary~\ref{B:poset-op} gives (St). Finite intersections
are objects by heredity, and the linear order makes every finite
setwise stabilizer pointwise. The minimal-support argument of
Proposition~\ref{B:intersection} now gives (Op). The proof deliberately
uses heredity and all induced embeddings; for strong subclasses their
replacement is the explicit admissible-bridge hypothesis of
Remark~\ref{B:bridge-scope}.
\end{proof}

\subsection{Recovery of the general Kan realization}

\begin{theorem}[Known and rechecked: DF A.13, including its open subgroups]
\label{B:kan-recovery}
Let $K$ be a nonempty connected small Kan complex and put
$A(K)=\Delta_{\mathrm{inj}}/K$, with every simplex, including
degenerate simplices, allowed as an object. Then
$BA(K)\simeq |K|$. Every component of every covariant action
on $A(K)$ is either weakly contractible or weakly equivalent,
over $|K|$, to a connected covering space of $|K|$.
Every connected cover occurs.

For countable $K$, this is a class in (S), with (Rig), (SA),
and (FA). If $K$ is simply connected it also has (St) and
(Op). For general countable connected $K$, let
$\Gamma=\pi_1|K|$. Its limit automorphism group $\widehat G$
has a continuous open quotient
\[
 1\longrightarrow G\longrightarrow\widehat G
       \xrightarrow{\rho}\Gamma\longrightarrow1,
\]
where $G$ is the automorphism group of the rich labelled
limit for the universal cover $\widetilde K$.
Every open subgroup of $\widehat G$ is either $G_F$ for a
nonempty finite subset $F$ of that limit, or
$\rho^{-1}(H)$ for a subgroup $H\leq\Gamma$.
The first kind has contractible $E_U$; the second has $E_U$
weakly equivalent to the connected cover corresponding to $H$.
Thus the theorem specifies all higher homotopy groups and their
monodromy, including the subgroups that are not finite
pointwise stabilizers.
\end{theorem}
\begin{proof}
We give the checked steps, since the passage from a simply
connected label complex to a general one is essential.

First, the simplex-category comparison used in
Proposition~\ref{B:cocycle-audit} works for any simplicial set:
the comma category at a simplex of $\Delta/K$ is the category
of its increasing-domain factorizations. Prepending a first
vertex with the first value of that simplex gives a functor
receiving natural maps from both the identity and a constant
one-vertex object. This contracts the comma category. Hence
Quillen's comparison gives
$B(\Delta_{\mathrm{inj}}/K)\simeq B(\Delta/K)\simeq |K|$,
naturally in $K$.

We use the following precise extension fact. If
$L\subseteq\Delta^U$ is a finite subcomplex and a map
$L\to K$ is nullhomotopic, it extends as an actual simplicial
map on $\Delta^U$. Indeed restriction of simplicial mapping
spaces is a Kan fibration. Its vertex image is a union of
components: lift an edge path from an extendible constant map
to the given map. The mapping-space assertion and its
pushout-product proof are Proposition I.5.2 and Corollary
I.5.3(2) of \cite{ABGoerssJardine}, opened and checked here.
For a domain with finitely many contractible components,
connectedness of $K$ lets us join their constant values first;
the same conclusion follows. For a domain homotopy equivalent
to a circle it follows if $K$ is simply connected.

For (FA), fix a choice of filler for every horn into $K$.
Given $P\subseteq X,Z$ with $P\ne\varnothing$, take the
tagged union, order side points first-side-before-second-side
in each gap of $P$, and put $p=\min P$. Retain the prescribed
labels on subsets contained in one side. For every mixed
$T$ not containing $p$, in increasing order of cardinality,
the already assigned labels on
$\{p\}\cup(T\setminus\{t\})$, $t\in T$, form the horn
on $\{p\}\cup T$ missing the face $T$. The chosen filler
labels the full simplex and its missing face. Induction on
$|T|$ verifies all face identities: their intersections have
already been assigned. An embedding in the variable side
over $P$ preserves the anchor, order, and mixed subsets, and
therefore encounters exactly the same horns and fillers.
This is the required functorial amalgam. It is strong because
the underlying union is disjoint off $P$.

Suppose next that $K$ is countable and simply connected.
The finite labelled simplices form a countable hereditary
class. Joint embedding follows from the extension fact on
two disjoint simplices; amalgamation follows from the horn
construction. Thus there is a countable rich homogeneous
limit $M$. Strong duplicate amalgams show that every point
outside a finite $B$ has an infinite orbit under $G_B$;
in particular $\Fix_M(G_B)=B$.

We check the stabilizer step rather than assuming elimination
of imaginaries. Use the finite coset-path construction in
Proposition~\ref{B:cocycle-audit}. Its two extension problems
have domains that are unions of simplices. To connect two
neighboring copies $B_0,B_1$ through $C'$, the three
prescribed pieces are the whole simplices on
$B_0\cup B_1$, $B_0\cup C'$, and $B_1\cup C'$.
The intersections have nerve a filled triangle when their
common intersection $D$ is nonempty, and a circle when
$D$ is empty. Their map to $K$ is nullhomotopic in either
case. The other extension, creating the chain of intermediate
copies, prescribes the existing whole simplex on
$B_0\cup B_1$ and the new copies; intersections are exactly
$D$, so the domain has contractible components. The extension
fact solves both problems literally. The independent order
is retained by the same finite gap interpolation. Every
finite coset path therefore exists, proving
\[
 \langle G_B,G_C\rangle=G_{B\cap C}.
\]
All finite subsets are admissible. Choosing a finite support
of least cardinality for an open subgroup and using this
identity shows that each of its elements preserves that
support setwise. The linear order then makes it pointwise
fixed. This proves (Op), including the empty support for
$G$ itself. Theorem~\ref{B:saturation} and
Theorem~\ref{B:op-total} now give the all-action description
in the simply connected countable case.

For a countable connected $K$, its simplicial universal cover
$\widetilde K$ and deck group $\Gamma$ are countable.
Take the rich limit $M$ for $\widetilde K$. Let
$\widehat G$ consist of increasing permutations $u$ of $M$
for which some $\gamma\in\Gamma$ sends the label of every
finite simplex $S$ to the label of $uS$. Call $\gamma$ the
degree of $u$. It is unique, since a deck transformation
fixing a vertex is the identity. Twisting all labels by a
fixed deck transformation gives another rich limit;
back-and-forth gives an element of every degree. Ordinary
homogeneity after undoing that degree proves extension of
every finite semilinear embedding. The degree map is an
open continuous surjection: a one-point stabilizer lies in
its kernel, and the target is discrete. Its kernel is $G$.

The projected $K$-labelled structure on $M$ is the rich limit
for $A(K)$. To see this, every simplex has a lift, and any
extension of a projected finite simplex has a unique lift
after one lifted vertex is fixed. It is then realized by
richness for $\widetilde K$. Every finite projected
embedding has a unique deck degree and extends as above.
For an arbitrary projected automorphism, the degrees on
overlapping finite sets agree on a vertex; taking unions
makes its degree global. Its full automorphism group is
therefore exactly $\widehat G$.

Let $U\leq\widehat G$ be open. The simply connected case
gives $U\cap G=G_F$ for a finite $F$, possibly empty.
Normality of $G$ gives $G_{uF}=G_F$ for every $u\in U$,
so the fixed-point identity gives $uF=F$. If $F$ is nonempty,
the independent linear order makes every $u$ fix it
pointwise. Its degree is then trivial, and $U=G_F$.
If $F$ is empty, $U$ contains $G$ and equals
$\rho^{-1}(H)$ for a subgroup $H\leq\Gamma$.

For nonempty $F$, the orbit action associated with $G_F$
is representable, so its element category is contractible.
For $U=\rho^{-1}(H)$ it is the locally constant action on
$\Gamma/H$ in which an embedding of degree $\gamma$
acts by left multiplication by $\gamma$. Its element
category is a covering category of $A(K)$.
For $H=1$, identify it with $A(\widetilde K)$ by sending
an object $(B,\gamma)$ to the simplex with label
$\gamma^{-1}\sigma_B$. If an arrow has degree $\delta$
and $\gamma'=\delta\gamma$, the transformed target label
restricts to
$\gamma'^{-1}\delta\sigma_B=\gamma^{-1}\sigma_B$.
This proves full faithfulness, and lifting of simplices
proves essential surjectivity. This cover is connected and
simply connected, so it is the universal cover. Quotienting
its regular fibre by $H$ identifies the general covering
with the specified monodromy. The saturation theorem now
proves the assertion for every action, not only orbit actions.

Finally let $K$ be any small connected Kan complex. Its
countable connected Kan subcomplexes form a filtered covering:
close a countable set of simplices under faces, degeneracies,
paths to one vertex, and chosen horn fillers in countably
many stages. The same construction applied to two such sets
proves filteredness. Fix a base object in an action component
$D$, and for each sufficiently large subcomplex $K_i$ let
$D_i$ be the component containing it in the restricted
element category. Finite zigzags show
\[
 ND=\bigcup_i ND_i.
\]
Homotopy groups commute with these filtered unions: after
the usual finite simplicial subdivisions, a representative
and a nullhomotopy involve only finitely many simplices and
occur at one stage. Each map $BD_i\to|K_i|$ is injective
on fundamental groups, so the union map is also injective:
a loop and a nullhomotopy of its image occur together at a
later stage.

If $BD$ is not weakly contractible, choose a nonzero homotopy
class in it. From one stage onward its representative
survives at every later stage. Each of these stages is
therefore of covering type, and its projection induces an
isomorphism on every homotopy group in degree at least two.
The union map has the same property. Lifting it to the
connected cover of $|K|$ determined by its fundamental-group
image gives a weak equivalence. Conversely any connected
simplicial cover $K_H\to K$ induces the discrete
opfibration $A(K_H)\to A(K)$: a lift of one face uniquely
extends across the whole simplex. Its associated action has
the required covering element category. This proves all
assertions for arbitrary small $K$.
\end{proof}

\begin{corollary}[Proved: every requested connectivity gap occurs]
\label{B:connectivity-gaps}
For each $n\geq2$ there is a rigid class in (S) and an open
subgroup $U$ whose $E_U$ is $(n-1)$-connected and not
$n$-connected. One may require $U$ not to be a finite
pointwise stabilizer. Moreover, for connected Kan $K$,
\[
 A(K)\in\mathcal Q_n\quad\Longleftrightarrow\quad
 |K|\text{ is }(n-1)\text{-connected}.
\]
\end{corollary}
\begin{proof}
The singleton action proves necessity in the displayed
equivalence. If the right side holds, $K$ is simply connected,
so every action component is contractible or weakly equivalent
to $K$; this proves sufficiency. The cocycle example
$K=K(\F_2,n)$ realizes the first assertion with $U=G$.
For the stronger requirement take the countable Kan complex
$K=K(\F_2,n)\times B(\Z/2)$ and let $U$ be the kernel
of the degree quotient in Theorem~\ref{B:kan-recovery}.
It is proper and open, and $E_U$ has the homotopy type of
the universal cover $K(\F_2,n)$. It cannot equal $G_F$
for nonempty finite $F$, since all those element categories
are contractible; nor is it the empty-set stabilizer, which
is the whole group.
\end{proof}
